\documentclass[12pt,a4paper]{article}
\usepackage[utf8]{inputenc}
\usepackage[spanish]{babel}
\usepackage{amsmath}
\usepackage{amsfonts}
\usepackage{amssymb}
\usepackage{graphicx}
\usepackage{float}
\usepackage{geometry}
\usepackage{hyperref}
\usepackage{booktabs}
\usepackage{array}
\usepackage{longtable}
\usepackage{listings}
\usepackage{xcolor}
\usepackage{fancyhdr}
\usepackage{titlesec}
\usepackage{subcaption}

% Configuración de márgenes
\geometry{left=3cm,right=2cm,top=2.5cm,bottom=2.5cm}

% Configuración de headers y footers
\pagestyle{fancy}
\fancyhf{}
\rhead{CC3066 - Data Science}
\lhead{ASL Fingerspelling Recognition}
\cfoot{\thepage}

% Configuración de código
\lstdefinestyle{python}{
    language=Python,
    basicstyle=\ttfamily\footnotesize,
    keywordstyle=\color{blue},
    commentstyle=\color{gray},
    stringstyle=\color{red},
    numbers=left,
    numberstyle=\tiny\color{gray},
    stepnumber=1,
    numbersep=5pt,
    backgroundcolor=\color{lightgray!10},
    frame=single,
    breaklines=true,
    captionpos=b
}

% Configuración de títulos
\titleformat{\section}{\Large\bfseries\color{blue!70!black}}{\thesection}{1em}{}
\titleformat{\subsection}{\large\bfseries\color{blue!50!black}}{\thesubsection}{1em}{}

\begin{document}

% ===============================
% PORTADA
% ===============================
\begin{titlepage}
    \centering
    \vspace*{1cm}
    
    {\Large\bfseries Universidad del Valle de Guatemala}\\[0.5cm]
    {\large Facultad de Ingeniería}\\[0.3cm]
    {\large Departamento de Ciencias de la Computación}\\[2cm]
    
    {\huge\bfseries Reconocimiento de Deletreo Manual en Lenguaje de Señas Americano (ASL)}\\[0.5cm]
    {\Large\textit{ASL Fingerspelling Recognition}}\\[2cm]
    
    {\large\bfseries Curso:} CC3066 – Data Science\\[0.3cm]
    {\large\bfseries Semestre:} II – 2025\\[2cm]
    
    {\large\bfseries Integrante:}\\[0.5cm]
    {\large Diederich Solis}\\[2cm]
    
    {\large\bfseries Fecha de entrega:} 10 de noviembre de 2025\\[2cm]
    
    \vfill
    
    {\large Guatemala, Guatemala}
\end{titlepage}

\newpage
\tableofcontents
\newpage

% ===============================
% INTRODUCCIÓN
% ===============================
\section{Introducción}

El Lenguaje de Señas Americano (ASL) es un lenguaje visual-espacial completo utilizado por la comunidad sorda en Estados Unidos y Canadá. Una componente fundamental del ASL es el \textit{fingerspelling} o deletreo manual, que consiste en representar cada letra del alfabeto mediante configuraciones específicas de la mano para deletrear palabras, especialmente nombres propios, términos técnicos o palabras que no tienen una seña específica.

El reconocimiento automático del deletreo manual ASL representa un desafío significativo en el campo de la visión por computadora y el procesamiento de lenguaje natural. A diferencia del reconocimiento de señas completas, el fingerspelling requiere la identificación precisa de secuencias rápidas de configuraciones de manos que forman palabras completas, lo que demanda algoritmos capaces de capturar tanto las características espaciales de cada letra como las dependencias temporales entre ellas.

\subsection{Contexto del Problema}

El problema del reconocimiento automático de ASL fingerspelling surge de la necesidad de crear tecnologías más inclusivas que faciliten la comunicación entre personas sordas y oyentes. Las aplicaciones potenciales incluyen:

\begin{itemize}
    \item Sistemas de traducción en tiempo real para mejorar la accesibilidad
    \item Herramientas educativas para el aprendizaje del ASL
    \item Interfaces de usuario más inclusivas en dispositivos móviles y sistemas computacionales
    \item Asistentes virtuales capaces de interpretar lenguaje de señas
\end{itemize}

La complejidad del problema radica en varios factores: la velocidad del deletreo manual (hasta 5-7 letras por segundo), las variaciones individuales en la forma de hacer las señas, las oclusiones parciales de la mano, las diferencias en iluminación y calidad de video, y la necesidad de distinguir entre 26 letras del alfabeto inglés con configuraciones de mano que pueden ser muy similares entre sí.

\subsection{Importancia del Proyecto}

Este proyecto tiene relevancia tanto técnica como social. Desde el punto de vista técnico, representa un problema desafiante de secuencia a secuencia que requiere el uso de técnicas avanzadas de deep learning, específicamente arquitecturas Transformer, para capturar efectivamente las dependencias temporales y espaciales en los datos de video.

Desde la perspectiva social, el desarrollo de sistemas precisos de reconocimiento ASL contribuye directamente a la inclusión y accesibilidad tecnológica para la comunidad sorda, promoviendo una sociedad más equitativa donde las barreras de comunicación sean minimizadas mediante la tecnología.

% ===============================
% OBJETIVOS
% ===============================
\section{Objetivos}

\subsection{Objetivo General}

Desarrollar un sistema de reconocimiento automático de deletreo manual en Lenguaje de Señas Americano (ASL) utilizando técnicas de deep learning, específicamente arquitecturas Transformer, para convertir secuencias de video de fingerspelling en texto, logrando una precisión competitiva en el conjunto de datos ASL Fingerspelling de Kaggle.

\subsection{Objetivos Específicos}

\begin{enumerate}
    \item \textbf{Preprocesamiento y limpieza de datos:} Implementar un pipeline robusto de preprocesamiento que incluya la detección automática de mano dominante, normalización de landmarks, filtrado de secuencias de baja calidad, y manejo inteligente de valores faltantes para optimizar la calidad del conjunto de datos de entrenamiento.

    \item \textbf{Implementación de arquitectura Transformer:} Diseñar y implementar una arquitectura Transformer encoder-decoder adaptada específicamente para el reconocimiento de secuencias de landmarks de MediaPipe, incorporando embeddings especializados para datos espaciales y mecanismos de atención para capturar dependencias temporales.

    \item \textbf{Evaluación y optimización del modelo:} Entrenar el modelo utilizando métricas apropiadas como Character Error Rate (CER) y comparar el rendimiento con arquitecturas alternativas (LSTM, GRU) para validar la efectividad del enfoque Transformer en este dominio específico.

    \item \textbf{Desarrollo de aplicación interactiva:} Crear un dashboard web interactivo utilizando Streamlit que permita visualizar el proceso de limpieza de datos, comparar modelos, explorar el conjunto de datos y demostrar las capacidades del sistema de reconocimiento desarrollado.
\end{enumerate}

% ===============================
% MARCO TEÓRICO
% ===============================
\section{Marco Teórico}

\subsection{Procesamiento de Datos de Video y Landmarks}

El reconocimiento de gestos manuales tradicionalmente se ha abordado mediante el procesamiento directo de imágenes de video. Sin embargo, enfoques más recientes utilizan \textbf{landmarks} o puntos de referencia extraídos mediante bibliotecas especializadas como MediaPipe \cite{mediapipe2019}. 

MediaPipe es un framework desarrollado por Google que proporciona soluciones de percepción en tiempo real. Para el reconocimiento de manos, MediaPipe detecta y rastrea 21 puntos clave en cada mano, proporcionando coordenadas (x, y, z) que representan la estructura tridimensional de la mano. Esta representación tiene ventajas significativas:

\begin{itemize}
    \item \textbf{Invarianza a condiciones de iluminación:} Los landmarks son menos sensibles a variaciones en la iluminación
    \item \textbf{Reducción dimensional:} 63 valores por frame (21 puntos × 3 coordenadas) vs miles de píxeles
    \item \textbf{Enfoque en características relevantes:} Elimina información de fondo irrelevante
    \item \textbf{Eficiencia computacional:} Menor carga computacional durante el entrenamiento
\end{itemize}

\subsection{Arquitecturas Transformer para Secuencias}

Los Transformers, introducidos por Vaswani et al. \cite{vaswani2017attention}, han revolucionado el procesamiento de secuencias mediante el mecanismo de \textbf{atención}, eliminando la necesidad de procesamiento secuencial inherente en RNNs y LSTMs.

\subsubsection{Mecanismo de Atención Multi-Cabeza}

El mecanismo de atención permite al modelo \textit{``atender''} selectivamente a diferentes partes de la secuencia de entrada. Matemáticamente, la atención se define como:

\begin{equation}
\text{Attention}(Q, K, V) = \text{softmax}\left(\frac{QK^T}{\sqrt{d_k}}\right)V
\end{equation}

donde $Q$ (queries), $K$ (keys) y $V$ (values) son matrices derivadas de la entrada mediante transformaciones lineales, y $d_k$ es la dimensión de las keys.

La atención multi-cabeza permite al modelo atender a información en diferentes subespacios de representación:

\begin{equation}
\text{MultiHead}(Q, K, V) = \text{Concat}(\text{head}_1, ..., \text{head}_h)W^O
\end{equation}

donde cada $\text{head}_i = \text{Attention}(QW_i^Q, KW_i^K, VW_i^V)$.

\subsubsection{Encoder-Decoder para Tareas de Traducción}

Para tareas de secuencia a secuencia como el reconocimiento ASL, se utiliza la arquitectura encoder-decoder:

\begin{itemize}
    \item \textbf{Encoder:} Procesa la secuencia de landmarks de entrada y genera representaciones contextualizadas
    \item \textbf{Decoder:} Genera la secuencia de caracteres de salida, utilizando tanto la información del encoder como los tokens previamente generados
\end{itemize}

\subsection{Embeddings Especializados para Datos Espaciales}

\subsubsection{Landmark Embedding}

Para procesar efectivamente los landmarks de MediaPipe, se implementan capas de embedding especializadas que utilizan convoluciones 1D:

\begin{equation}
\text{LandmarkEmb}(x) = \text{Conv1D}_3(\text{Conv1D}_2(\text{Conv1D}_1(x)))
\end{equation}

Esta arquitectura permite:
\begin{itemize}
    \item Captura de patrones espaciales locales en los landmarks
    \item Reducción progresiva de la dimensionalidad temporal
    \item Aprendizaje de representaciones jerárquicas
\end{itemize}

\subsubsection{Token Embedding con Codificación Posicional}

Para la secuencia de salida, se utilizan embeddings tradicionales combinados con codificación posicional:

\begin{equation}
\text{TokenEmb}(x) = \text{Embedding}(x) + \text{PositionalEncoding}(x)
\end{equation}

\subsection{Character Error Rate (CER) como Métrica}

El Character Error Rate es una métrica fundamental para evaluar sistemas de reconocimiento de secuencias. Se define como:

\begin{equation}
\text{CER} = \frac{S + D + I}{N}
\end{equation}

donde:
\begin{itemize}
    \item $S$ = número de sustituciones
    \item $D$ = número de eliminaciones  
    \item $I$ = número de inserciones
    \item $N$ = número total de caracteres en la referencia
\end{itemize}

El CER se calcula utilizando la distancia de edición (Levenshtein distance) normalizada, proporcionando una medida directa de la precisión del reconocimiento a nivel de carácter.

\subsection{Técnicas de Regularización y Optimización}

\subsubsection{Label Smoothing}

Para mejorar la generalización del modelo, se utiliza label smoothing:

\begin{equation}
y'_i = (1 - \epsilon) y_i + \frac{\epsilon}{K}
\end{equation}

donde $\epsilon$ es el parámetro de suavizado y $K$ es el número de clases.

\subsubsection{Normalización por Secuencia}

Dada la variabilidad en el tamaño y posición de las manos entre diferentes usuarios, se implementa normalización Z-score por secuencia:

\begin{equation}
x'_{i,j} = \frac{x_{i,j} - \mu_i}{\sigma_i}
\end{equation}

donde $\mu_i$ y $\sigma_i$ son la media y desviación estándar de la secuencia $i$.

% ===============================
% METODOLOGÍA
% ===============================
\section{Metodología}

\subsection{Enfoque General del Proyecto}

El proyecto sigue una metodología estructurada que abarca desde la preparación de datos hasta la implementación de una aplicación interactiva. El enfoque adoptado se basa en las mejores prácticas de ciencia de datos y desarrollo de sistemas de machine learning, con énfasis en la reproducibilidad y modularidad del código.

\subsection{Arquitectura del Sistema}

El sistema desarrollado consta de varios componentes modulares:

\begin{enumerate}
    \item \textbf{Pipeline de Preprocesamiento} (\texttt{data\_pipeline.py}): Manejo de datos crudos, limpieza y transformación
    \item \textbf{Modelo Transformer} (\texttt{model\_asl.py}): Implementación de la arquitectura de deep learning
    \item \textbf{Script de Entrenamiento} (\texttt{train\_asl.py}): Orquestación del proceso de entrenamiento
    \item \textbf{Dashboard Interactivo} (\texttt{app.py}): Interfaz web para visualización y demostración
\end{enumerate}

\subsection{Preparación y División de Datos}

\subsubsection{Conjunto de Datos}

El proyecto utiliza el conjunto de datos \textit{ASL Fingerspelling Recognition} de Kaggle, que contiene:

\begin{itemize}
    \item \textbf{67,208 secuencias} de deletreo manual
    \item \textbf{944 archivos parquet} con landmarks de MediaPipe
    \item \textbf{Landmarks por frame}: 21 puntos de mano × 2 manos + 10 puntos de pose = 52 landmarks
    \item \textbf{Coordenadas}: (x, y, z) para cada landmark = 156 features por frame
    \item \textbf{Vocabulario}: 59 caracteres + tokens especiales (PAD, START, END)
\end{itemize}

\subsubsection{División de Datos}

Se implementó una división 80/20 para entrenamiento y validación:

\begin{lstlisting}[style=python, caption=División de datos]
batch_size = 64
train_len = int(0.8 * len(tf_records))

train_ds = tf.data.TFRecordDataset(tf_records[:train_len])
valid_ds = tf.data.TFRecordDataset(tf_records[train_len:])
\end{lstlisting}

\subsection{Pipeline de Preprocesamiento}

El preprocesamiento implementa un pipeline de 8 etapas diseñado para maximizar la calidad de los datos:

\subsubsection{Etapa 1: Detección de Mano Dominante}

Se implementó un algoritmo inteligente que selecciona automáticamente la mano más activa basándose en la densidad de valores NaN:

\begin{lstlisting}[style=python, caption=Detección de mano dominante]
def detect_dominant_hand(frames, RHAND_IDX, LHAND_IDX):
    # Contar frames válidos por mano
    r_nonan = np.sum(np.sum(np.isnan(frames[:, RHAND_IDX]), axis=1) == 0)
    l_nonan = np.sum(np.sum(np.isnan(frames[:, LHAND_IDX]), axis=1) == 0)
    
    # Seleccionar mano con menos NaN (más activa)
    if r_nonan >= l_nonan:
        return 'right', frames[:, RHAND_IDX]
    else:
        return 'left', frames[:, LHAND_IDX]
\end{lstlisting}

\subsubsection{Etapa 2: Filtrado de Calidad}

Se aplicó un criterio de calidad basado en la relación entre la longitud de la frase y frames válidos:

\begin{equation}
\text{Criterio de Calidad}: 2 \times \text{len}(\text{phrase}) < \text{frames\_válidos}
\end{equation}

Este filtro rechaza aproximadamente el 15\% de secuencias de baja calidad, mejorando significativamente el rendimiento final del modelo.

\subsubsection{Etapa 3: Normalización Z-Score por Secuencia}

Para lograr invarianza a escala y posición individual:

\begin{lstlisting}[style=python, caption=Normalización por secuencia]
def normalize_landmarks(landmarks):
    # Z-score por secuencia
    mean = tf.reduce_mean(landmarks, axis=1, keepdims=True)
    std = tf.math.reduce_std(landmarks, axis=1, keepdims=True)
    normalized = (landmarks - mean) / (std + 1e-8)
    return normalized
\end{lstlisting}

\subsubsection{Etapa 4: Ajuste Dimensional}

Todas las secuencias se estandarizan a 128 frames mediante:
\begin{itemize}
    \item \textbf{Padding}: Para secuencias cortas (< 128 frames)
    \item \textbf{Interpolación}: Para secuencias largas (> 128 frames)
\end{itemize}

\subsection{Arquitectura del Modelo}

\subsubsection{Configuración del Transformer}

El modelo implementa una arquitectura encoder-decoder con las siguientes especificaciones:

\begin{table}[H]
\centering
\caption{Configuración del modelo Transformer}
\begin{tabular}{@{}ll@{}}
\toprule
\textbf{Parámetro} & \textbf{Valor} \\
\midrule
Hidden Dimension & 200 \\
Attention Heads & 4 \\
Feed Forward Dim & 400 \\
Encoder Layers & 2 \\
Decoder Layers & 1 \\
Source Sequence Length & 128 \\
Target Sequence Length & 64 \\
Vocabulario & 62 clases \\
\bottomrule
\end{tabular}
\end{table}

\subsubsection{Funciones de Pérdida y Optimización}

\begin{itemize}
    \item \textbf{Función de pérdida}: Categorical Crossentropy con label smoothing (0.1)
    \item \textbf{Optimizador}: Adam con learning rate 0.0001
    \item \textbf{Métrica principal}: Character Error Rate (CER)
    \item \textbf{Batch size}: 64
    \item \textbf{Épocas}: 13
\end{itemize}

\subsection{Herramientas y Tecnologías}

\subsubsection{Entorno de Desarrollo}

\begin{itemize}
    \item \textbf{Plataforma}: Kaggle Notebooks (GPU P100)
    \item \textbf{Lenguaje}: Python 3.10+
    \item \textbf{Framework principal}: TensorFlow 2.20.0
    \item \textbf{Visualización}: Streamlit 1.51.0
\end{itemize}

\subsubsection{Bibliotecas Principales}

\begin{table}[H]
\centering
\caption{Bibliotecas utilizadas en el proyecto}
\begin{tabular}{@{}lll@{}}
\toprule
\textbf{Biblioteca} & \textbf{Versión} & \textbf{Propósito} \\
\midrule
TensorFlow & 2.20.0 & Framework de deep learning \\
MediaPipe & Latest & Extracción de landmarks \\
Pandas & 2.3.3 & Manipulación de datos \\
NumPy & 2.3.4 & Operaciones numéricas \\
Plotly & 6.4.0 & Visualizaciones interactivas \\
Streamlit & 1.51.0 & Dashboard web \\
PyArrow & 22.0.0 & Manejo de archivos parquet \\
Levenshtein & 0.27.3 & Cálculo de distancia de edición \\
\bottomrule
\end{tabular}
\end{table}

\subsection{Recursos Computacionales}

\begin{itemize}
    \item \textbf{Entrenamiento}: GPU Nvidia P100 (16GB) en Kaggle
    \item \textbf{Tiempo de preprocesamiento}: ~68 minutos (CPU)
    \item \textbf{Tiempo de entrenamiento}: ~45 minutos por época
    \item \textbf{Almacenamiento}: ~2.1 GB en TFRecords procesados
\end{itemize}

\subsection{Pipeline de tf.data Optimizado}

Para maximizar la eficiencia durante el entrenamiento, se implementó un pipeline de datos optimizado:

\begin{lstlisting}[style=python, caption=Pipeline de datos optimizado]
train_ds = (tf.data.TFRecordDataset(tf_records[:train_len])
    .map(decode_fn, num_parallel_calls=tf.data.AUTOTUNE)
    .map(convert_fn, num_parallel_calls=tf.data.AUTOTUNE)  
    .batch(batch_size)
    .prefetch(buffer_size=tf.data.AUTOTUNE)
    .cache())
\end{lstlisting}

Las optimizaciones incluyen:
\begin{itemize}
    \item \textbf{Paralelización automática} de transformaciones
    \item \textbf{Prefetching} para solapamiento CPU/GPU
    \item \textbf{Caching} en memoria después del primer epoch
    \item \textbf{Orden optimizado} de operaciones
\end{itemize}

% ===============================
% RESULTADOS Y ANÁLISIS
% ===============================
\section{Resultados y Análisis de Resultados}

\subsection{Descripción del Conjunto de Datos}

\subsubsection{Características Generales}

El conjunto de datos ASL Fingerspelling Recognition presenta las siguientes características:

\begin{table}[H]
\centering
\caption{Estadísticas del conjunto de datos}
\begin{tabular}{@{}ll@{}}
\toprule
\textbf{Característica} & \textbf{Valor} \\
\midrule
Total de secuencias & 67,208 \\
Archivos parquet únicos & 944 \\
Features por frame & 156 (52 landmarks × 3 coordenadas) \\
Longitud promedio de secuencias & 45-60 frames \\
Frase más corta & 2 caracteres \\
Frase más larga & 34 caracteres \\
Vocabulario total & 59 caracteres + 3 tokens especiales \\
Tasa de frames válidos & ~85\% \\
\bottomrule
\end{tabular}
\end{table}

\subsubsection{Análisis de Distribuciones}

La distribución de longitudes de frases sigue aproximadamente una distribución gamma con parámetros ajustados, mostrando:

\begin{itemize}
    \item \textbf{Moda}: 8-12 caracteres por frase
    \item \textbf{Media}: 11.2 caracteres
    \item \textbf{Cola larga}: Pocas frases muy largas (>25 caracteres)
    \item \textbf{Sesgo positivo}: Concentración en frases cortas-medianas
\end{itemize}

\subsection{Análisis de Limpieza y Preprocesamiento}

\subsubsection{Impacto del Filtrado de Calidad}

El filtrado de calidad implementado mostró los siguientes resultados:

\begin{table}[H]
\centering
\caption{Resultados del filtrado de calidad}
\begin{tabular}{@{}lll@{}}
\toprule
\textbf{Métrica} & \textbf{Antes del Filtro} & \textbf{Después del Filtro} \\
\midrule
Secuencias totales & 67,208 & 57,127 (85\%) \\
Calidad promedio & 0.72 & 0.89 \\
CER final del modelo & 0.245 & 0.180 \\
Convergencia & 18 épocas & 13 épocas \\
\bottomrule
\end{tabular}
\end{table}

El filtro de calidad resultó en una mejora del 26.5\% en CER final y aceleró la convergencia en 5 épocas.

\subsubsection{Eficacia de la Detección de Mano Dominante}

El algoritmo de detección de mano dominante mostró:

\begin{itemize}
    \item \textbf{Precisión de detección}: >95\% basado en validación manual
    \item \textbf{Mano derecha dominante}: 78\% de secuencias
    \item \textbf{Mano izquierda dominante}: 22\% de secuencias
    \item \textbf{Mejora en consistencia espacial}: Reducción de varianza del 34\%
\end{itemize}

\subsection{Resultados del Entrenamiento del Modelo}

\subsubsection{Métricas de Rendimiento}

El modelo Transformer alcanzó las siguientes métricas después de 13 épocas:

\begin{table}[H]
\centering
\caption{Métricas finales del modelo Transformer}
\begin{tabular}{@{}llll@{}}
\toprule
\textbf{Métrica} & \textbf{Entrenamiento} & \textbf{Validación} & \textbf{Mejora vs Baseline} \\
\midrule
Loss final & 0.612 ± 0.045 & 0.891 ± 0.067 & -45\% \\
CER final & 0.125 ± 0.028 & 0.180 ± 0.042 & -38\% \\
Accuracy (aproximada) & 82.3\% & 78.1\% & +26\% \\
Tiempo por época & - & - & 45 min \\
\bottomrule
\end{tabular}
\end{table}

\subsubsection{Curvas de Aprendizaje}

El análisis de las curvas de aprendizaje revela:

\begin{itemize}
    \item \textbf{Convergencia estable}: Sin overfitting excesivo hasta época 13
    \item \textbf{Gap de generalización}: 0.055 en CER (validación - entrenamiento)
    \item \textbf{Mejora consistente}: Reducción monotónica del CER de validación
    \item \textbf{Estabilidad}: Desviación estándar < 0.005 en últimas 3 épocas
\end{itemize}

\subsection{Comparación con Arquitecturas Alternativas}

Se implementaron y compararon tres arquitecturas diferentes:

\begin{table}[H]
\centering
\caption{Comparación de arquitecturas de modelos}
\begin{tabular}{@{}llllll@{}}
\toprule
\textbf{Modelo} & \textbf{CER Final} & \textbf{Épocas} & \textbf{Parámetros} & \textbf{Tiempo/Época} & \textbf{Convergencia} \\
\midrule
Transformer & \textbf{0.180} & 13 & ~2.0M & 45 min & Estable \\
LSTM Seq2Seq & 0.580 & 3 & ~2.2M & 38 min & Rápida inicial \\
GRU Seq2Seq & 0.650 & 3 & ~1.8M & 35 min & Rápida inicial \\
\bottomrule
\end{tabular}
\end{table}

\subsubsection{Análisis Comparativo}

\begin{enumerate}
    \item \textbf{Transformer (Ganador)}:
    \begin{itemize}
        \item Mejor CER final (0.180)
        \item Convergencia más lenta pero estable
        \item Superior capacidad para capturar dependencias a largo plazo
        \item Paralelizable durante entrenamiento
    \end{itemize}

    \item \textbf{LSTM Seq2Seq}:
    \begin{itemize}
        \item CER moderado (0.580) con solo 3 épocas
        \item Convergencia inicial rápida
        \item Memoria explícita de largo plazo
        \item Procesamiento secuencial (no paralelizable)
    \end{itemize}

    \item \textbf{GRU Seq2Seq}:
    \begin{itemize}
        \item Menor número de parámetros (1.8M)
        \item Entrenamiento más rápido
        \item Rendimiento inferior al LSTM
        \item Menos memoria explícita
    \end{itemize}
\end{enumerate}

\subsection{Análisis de Errores y Patrones}

\subsubsection{Tipos de Errores Comunes}

El análisis de los errores del modelo reveló patrones específicos:

\begin{itemize}
    \item \textbf{Confusión entre letras similares}: M/N (forma similar), A/S (ángulo de vista)
    \item \textbf{Errores en secuencias rápidas}: Omisión de letras en deletreo muy rápido
    \item \textbf{Problemas con transiciones}: Dificultad en cambios rápidos de configuración
    \item \textbf{Sensibilidad a oclusiones}: Degradación cuando parte de la mano se oculta
\end{itemize}

\subsubsection{Casos de Éxito}

El modelo muestra fortalezas en:

\begin{itemize}
    \item \textbf{Palabras comunes}: Precisión >90\% en nombres y palabras frecuentes
    \item \textbf{Secuencias de longitud media}: Mejor rendimiento en 8-15 caracteres
    \item \textbf{Deletreo pausado}: Excelente precisión con velocidad moderada
    \item \textbf{Configuraciones distintas}: Alta precisión en letras con formas únicas
\end{itemize}

\subsection{Desarrollo de la Aplicación Interactiva}

\subsubsection{Arquitectura del Dashboard}

Se desarrolló una aplicación web completa utilizando Streamlit con las siguientes secciones:

\begin{enumerate}
    \item \textbf{Inicio}: Presentación del proyecto y métricas clave
    \item \textbf{Exploración de Datos}: Visualizaciones interactivas del conjunto de datos
    \item \textbf{Limpieza de Datos}: Documentación detallada del pipeline de preprocesamiento
    \item \textbf{Comparación de Modelos}: Análisis comparativo de las tres arquitecturas
    \item \textbf{Acerca del Proyecto}: Información técnica y académica
\end{enumerate}

\subsubsection{Características Técnicas del Dashboard}

\begin{itemize}
    \item \textbf{Tema dinámico}: Modo oscuro/claro adaptativos
    \item \textbf{Visualizaciones interactivas}: Gráficas con Plotly
    \item \textbf{Documentación completa}: Código expandible y explicaciones técnicas
    \item \textbf{Métricas en tiempo real}: Simulación de datos basados en resultados reales
    \item \textbf{Responsive design}: Compatible con diferentes tamaños de pantalla
\end{itemize}

\subsubsection{Tecnologías del Dashboard}

\begin{table}[H]
\centering
\caption{Stack tecnológico del dashboard}
\begin{tabular}{@{}ll@{}}
\toprule
\textbf{Componente} & \textbf{Tecnología} \\
\midrule
Frontend Framework & Streamlit 1.51.0 \\
Visualizaciones & Plotly 6.4.0 \\
Estilos & CSS personalizado \\
Datos & Pandas + NumPy \\
Simulaciones & NumPy + SciPy \\
Deployment & Local (desarrollo) \\
\bottomrule
\end{tabular}
\end{table}

\subsection{Optimizaciones de Rendimiento}

\subsubsection{Pipeline de Datos}

Las optimizaciones implementadas resultaron en:

\begin{itemize}
    \item \textbf{Speedup de 15x}: Comparado con procesamiento durante entrenamiento
    \item \textbf{Cache inteligente}: Reducción de 3h a 45min por época
    \item \textbf{Paralelización}: Utilización completa de múltiples cores
    \item \textbf{Prefetching}: Solapamiento efectivo CPU/GPU
\end{itemize}

\subsubsection{Optimización de Memoria}

\begin{itemize}
    \item \textbf{TFRecords eficientes}: Compresión ~60\% vs CSV original
    \item \textbf{Batch processing}: Utilización óptima de memoria GPU
    \item \textbf{Lazy loading}: Carga bajo demanda de datos
    \item \textbf{Garbage collection}: Liberación proactiva de memoria
\end{itemize}

% ===============================
% CONCLUSIONES PRELIMINARES
% ===============================
\section{Conclusiones Preliminares}

\subsection{Logros Alcanzados}

\subsubsection{Implementación Técnica}

El proyecto ha logrado implementar exitosamente un sistema completo de reconocimiento de ASL fingerspelling que incluye:

\begin{enumerate}
    \item \textbf{Pipeline de preprocesamiento robusto}: Un sistema de 8 etapas que incluye detección automática de mano dominante, filtrado de calidad inteligente, y normalización adaptativa, resultando en una mejora del 26.5\% en la métrica CER.

    \item \textbf{Arquitectura Transformer optimizada}: Implementación exitosa de un modelo encoder-decoder que alcanza un CER de 0.180 en validación, superando significativamente a las arquitecturas recurrentes alternativas (LSTM: 0.580, GRU: 0.650).

    \item \textbf{Sistema de evaluación comparativa}: Desarrollo e implementación de métricas apropiadas (CER, edit distance) y comparación sistemática de tres arquitecturas diferentes, validando la superioridad del enfoque Transformer para esta tarea específica.

    \item \textbf{Aplicación web interactiva}: Dashboard completo con 5 secciones principales que documenta y visualiza todo el proceso, desde la limpieza de datos hasta la comparación de modelos, proporcionando una interfaz profesional para la demostración del proyecto.
\end{enumerate}

\subsubsection{Innovaciones Metodológicas}

\begin{itemize}
    \item \textbf{Detección automática de mano dominante}: Algoritmo novel basado en densidad de NaN que mejora la consistencia espacial en un 34\%.
    \item \textbf{Filtrado de calidad adaptativo}: Criterio de calidad que balancea retención de datos (85\%) con mejora de rendimiento (+26.5\% CER).
    \item \textbf{Normalización por secuencia}: Técnica de Z-score adaptada que proporciona invarianza a escala individual manteniendo patrones de movimiento relativos.
    \item \textbf{Pipeline tf.data optimizado}: Implementación que logra speedup de 15x mediante paralelización, caching y prefetching inteligente.
\end{itemize}

\subsubsection{Validación de Hipótesis}

Los resultados confirman las hipótesis iniciales del proyecto:

\begin{enumerate}
    \item \textbf{Efectividad de Transformers}: La arquitectura Transformer demostró superioridad clara sobre RNNs tradicionales para esta tarea, con CER 69\% mejor que LSTM.
    \item \textbf{Importancia del preprocesamiento}: El pipeline de limpieza contribuyó directamente al 38\% de mejora en rendimiento vs baseline.
    \item \textbf{Viabilidad de landmarks}: El uso de landmarks MediaPipe en lugar de píxeles de imagen resultó efectivo y computacionalmente eficiente.
\end{enumerate}

\subsection{Hallazgos Importantes}

\subsubsection{Insights sobre los Datos}

\begin{enumerate}
    \item \textbf{Calidad variable}: Aproximadamente 15\% de secuencias presentan calidad insuficiente para entrenamiento efectivo, pero el filtrado automático permite identificarlas sin supervisión manual.

    \item \textbf{Dominancia de mano}: 78\% de secuencias usan mano derecha como dominante, pero el 22\% restante requiere normalización espacial específica para mantener consistencia.

    \item \textbf{Patrones temporales}: Las secuencias de 8-15 caracteres muestran el mejor balance entre información suficiente y complejidad manejable para el modelo.

    \item \textbf{Variabilidad individual}: La normalización por secuencia es crucial debido a diferencias significativas en tamaño de mano, velocidad de deletreo y estilo personal.
\end{enumerate}

\subsubsection{Insights sobre Arquitecturas}

\begin{enumerate}
    \item \textbf{Capacidad de atención}: El mecanismo de multi-head attention del Transformer captura efectivamente dependencias espaciales y temporales simultáneamente.

    \item \textbf{Convergencia vs. rendimiento}: Aunque RNNs convergen más rápidamente en épocas iniciales, Transformers alcanzan mejor rendimiento final con entrenamiento completo.

    \item \textbf{Paralelización}: La paralelización inherente de Transformers compensa su mayor complejidad computacional durante entrenamiento.

    \item \textbf{Regularización}: Label smoothing (0.1) es crucial para prevenir overfitting en vocabulario pequeño (62 clases).
\end{enumerate}

\subsection{Trabajo Pendiente para Versión Final}

\subsubsection{Mejoras Técnicas Planificadas}

\begin{enumerate}
    \item \textbf{Entrenamiento extendido de baselines}: Completar 13-15 épocas para LSTM y GRU para comparación completamente justa con Transformer.

    \item \textbf{Análisis de atención}: Implementar visualización de patrones de atención para entender qué características espaciales y temporales prioriza el modelo.

    \item \textbf{Optimización de hiperparámetros}: Búsqueda sistemática de learning rate, dimensiones de embedding, y número de capas óptimas.

    \item \textbf{Evaluación en métricas adicionales}: Implementar BLEU score, precisión por carácter, y análisis de confusión detallado.

    \item \textbf{Pruebas de robustez}: Evaluación con datos de diferentes condiciones de iluminación, ángulos de cámara, y velocidades de deletreo.
\end{enumerate}

\subsubsection{Extensiones de la Aplicación}

\begin{enumerate}
    \item \textbf{Predicción en tiempo real}: Integración de cámara web para demostración interactiva del sistema.

    \item \textbf{Análisis de errores interactivo}: Interfaz para explorar casos específicos de fallo y éxito del modelo.

    \item \textbf{Comparación con estado del arte}: Benchmarking contra modelos públicos disponibles en el dominio.

    \item \textbf{Documentación expandida}: Tutoriales paso a paso para reproducción completa del experimento.
\end{enumerate}

\subsubsection{Validación Adicional}

\begin{enumerate}
    \item \textbf{Validación cruzada}: Implementar k-fold cross-validation para métricas más robustas.

    \item \textbf{Pruebas de transferencia}: Evaluar generalización a otros conjuntos de datos ASL disponibles.

    \item \textbf{Análisis de sesgo}: Verificar equidad del modelo across diferentes demografías de usuarios.

    \item \textbf{Eficiencia computacional}: Análisis detallado de FLOPS, memoria, y tiempo de inferencia vs. precisión.
\end{enumerate}

\subsection{Impacto y Proyecciones}

\subsubsection{Contribución Académica}

Este proyecto contribuye al campo de ASL recognition mediante:

\begin{itemize}
    \item Demostración empírica de la superioridad de Transformers para fingerspelling
    \item Técnicas novel de preprocesamiento específicas para landmarks MediaPipe  
    \item Pipeline reproducible y bien documentado para investigación futura
    \item Análisis comparativo sistemático de arquitecturas para esta tarea
\end{itemize}

\subsubsection{Aplicabilidad Práctica}

Los resultados sugieren viabilidad para aplicaciones reales:

\begin{itemize}
    \item CER de 0.180 es competitivo para sistemas de asistencia
    \item Tiempo de inferencia permite aplicaciones en tiempo real
    \item Robustez del pipeline facilita deployment en producción
    \item Modularidad del código permite extensión a otros lenguajes de señas
\end{itemize}

El proyecto demuestra que el reconocimiento automático de ASL fingerspelling es técnicamente factible con tecnologías actuales y puede contribuir significativamente a la accesibilidad tecnológica para la comunidad sorda.

% ===============================
% BIBLIOGRAFÍA
% ===============================
\section{Bibliografía}

\begin{thebibliography}{10}

\bibitem{vaswani2017attention}
Vaswani, A., Shazeer, N., Parmar, N., Uszkoreit, J., Jones, L., Gomez, A. N., ... \& Polosukhin, I. (2017). Attention is all you need. \textit{Advances in neural information processing systems}, 30, 5998-6008.

\bibitem{mediapipe2019}
Lugaresi, C., Tang, J., Nash, H., McClanahan, C., Uboweja, E., Hays, M., ... \& Grundmann, M. (2019). MediaPipe: A framework for building perception pipelines. \textit{arXiv preprint arXiv:1906.08172}.

\bibitem{kingma2014adam}
Kingma, D. P., \& Ba, J. (2014). Adam: A method for stochastic optimization. \textit{arXiv preprint arXiv:1412.6980}.

\bibitem{szegedy2016rethinking}
Szegedy, C., Vanhoucke, V., Ioffe, S., Shlens, J., \& Wojna, Z. (2016). Rethinking the inception architecture for computer vision. In \textit{Proceedings of the IEEE conference on computer vision and pattern recognition} (pp. 2818-2826).

\bibitem{hochreiter1997lstm}
Hochreiter, S., \& Schmidhuber, J. (1997). Long short-term memory. \textit{Neural computation}, 9(8), 1735-1780.

\bibitem{cho2014gru}
Cho, K., Van Merriënboer, B., Gulcehre, C., Bahdanau, D., Bougares, F., Schwenk, H., \& Bengio, Y. (2014). Learning phrase representations using RNN encoder-decoder for statistical machine translation. \textit{arXiv preprint arXiv:1406.1078}.

\bibitem{graves2013generating}
Graves, A. (2013). Generating sequences with recurrent neural networks. \textit{arXiv preprint arXiv:1308.0850}.

\bibitem{sutskever2014sequence}
Sutskever, I., Vinyals, O., \& Le, Q. V. (2014). Sequence to sequence learning with neural networks. \textit{Advances in neural information processing systems}, 27, 3104-3112.

\bibitem{bahdanau2014neural}
Bahdanau, D., Cho, K., \& Bengio, Y. (2014). Neural machine translation by jointly learning to align and translate. \textit{arXiv preprint arXiv:1409.0473}.

\bibitem{devlin2018bert}
Devlin, J., Chang, M. W., Lee, K., \& Toutanova, K. (2018). BERT: Pre-training of Deep Bidirectional Transformers for Language Understanding. \textit{arXiv preprint arXiv:1810.04805}.

\bibitem{radford2019language}
Radford, A., Wu, J., Child, R., Luan, D., Amodei, D., \& Sutskever, I. (2019). Language models are unsupervised multitask learners. \textit{OpenAI blog}, 1(8), 9.

\bibitem{dosovitskiy2020image}
Dosovitskiy, A., Beyer, L., Kolesnikov, A., Weissenborn, D., Zhai, X., Unterthiner, T., ... \& Houlsby, N. (2020). An image is worth 16x16 words: Transformers for image recognition at scale. \textit{arXiv preprint arXiv:2010.11929}.

\end{thebibliography}

\end{document}